\documentclass[11pt]{scrreprt}
%%fakesection Load packages
% These are evan.sty
\usepackage{amsmath,amssymb,amsthm}
\usepackage{mathrsfs}
\usepackage[usenames,svgnames,dvipsnames]{xcolor}
\usepackage{hyperref}
\usepackage[nameinlink]{cleveref}
\usepackage{textcomp}
\usepackage{enumerate}
\usepackage[obeyFinal,textsize=scriptsize,shadow]{todonotes}
\usepackage{mathtools}
\usepackage{microtype}

\usepackage[normalem]{ulem}
\usepackage{stmaryrd}
\usepackage{wasysym}
\usepackage{multirow}
\usepackage{prerex}

%%fakesection evan.sty macros
%Small commands
\newcommand{\cbrt}[1]{\sqrt[3]{#1}}
\newcommand{\floor}[1]{\left\lfloor #1 \right\rfloor}
\newcommand{\ceiling}[1]{\left\lceil #1 \right\rceil}
\newcommand{\mailto}[1]{\href{mailto:#1}{\texttt{#1}}}
\newcommand{\ol}{\overline}
\newcommand{\ul}{\underline}
\newcommand{\wt}{\widetilde}
\newcommand{\wh}{\widehat}
\newcommand{\eps}{\varepsilon}
%\renewcommand{\iff}{\Leftrightarrow}
%\renewcommand{\implies}{\Rightarrow}
\newcommand{\vocab}[1]{\textbf{\color{blue} #1}}
\newcommand{\catname}{\mathsf}
\newcommand{\hrulebar}{
  \par\hspace{\fill}\rule{0.95\linewidth}{.7pt}\hspace{\fill}
  \par\nointerlineskip \vspace{\baselineskip}
}
\newcommand{\half}{\frac{1}{2}}

%For use in author command
\newcommand{\plusemail}[1]{\\ \normalfont \texttt{\mailto{#1}}}

%More commands and math operators
\DeclareMathOperator{\cis}{cis}
\DeclareMathOperator*{\lcm}{lcm}
\DeclareMathOperator*{\argmin}{arg min}
\DeclareMathOperator*{\argmax}{arg max}

%Convenient Environments
\newenvironment{soln}{\begin{proof}[Solution]}{\end{proof}}
\newenvironment{parlist}{\begin{inparaenum}[(i)]}{\end{inparaenum}}
\newenvironment{gobble}{\setbox\z@\vbox\bgroup}{\egroup}

%Inequalities
\newcommand{\cycsum}{\sum_{\mathrm{cyc}}}
\newcommand{\symsum}{\sum_{\mathrm{sym}}}
\newcommand{\cycprod}{\prod_{\mathrm{cyc}}}
\newcommand{\symprod}{\prod_{\mathrm{sym}}}

%From H113 "Introduction to Abstract Algebra" at UC Berkeley
\newcommand{\CC}{\mathbb C}
\newcommand{\FF}{\mathbb F}
\newcommand{\NN}{\mathbb N}
\newcommand{\QQ}{\mathbb Q}
\newcommand{\RR}{\mathbb R}
\newcommand{\ZZ}{\mathbb Z}
\newcommand{\charin}{\text{ char }}
\DeclareMathOperator{\sign}{sign}
\DeclareMathOperator{\Aut}{Aut}
\DeclareMathOperator{\Inn}{Inn}
\DeclareMathOperator{\Syl}{Syl}
\DeclareMathOperator{\Gal}{Gal}
\DeclareMathOperator{\GL}{GL} % General linear group
\DeclareMathOperator{\SL}{SL} % Special linear group

%From Kiran Kedlaya's "Geometry Unbound"
\newcommand{\abs}[1]{\left\lvert #1 \right\rvert}
\newcommand{\norm}[1]{\left\lVert #1 \right\rVert}
\newcommand{\dang}{\measuredangle} %% Directed angle
\newcommand{\ray}[1]{\overrightarrow{#1}} 
\newcommand{\seg}[1]{\overline{#1}}
\newcommand{\arc}[1]{\wideparen{#1}}

%From M275 "Topology" at SJSU
\newcommand{\id}{\mathrm{id}}
\newcommand{\taking}[1]{\xrightarrow{#1}}
\newcommand{\inv}{^{-1}}

%From M170 "Introduction to Graph Theory" at SJSU
\DeclareMathOperator{\diam}{diam}
\DeclareMathOperator{\ord}{ord}
\newcommand{\defeq}{\overset{\mathrm{def}}{=}}

%From the USAMO .tex filse
\newcommand{\st}{^{\text{st}}}
\newcommand{\nd}{^{\text{nd}}}
\renewcommand{\th}{^{\text{th}}}
\newcommand{\dg}{^\circ}
\newcommand{\ii}{\item}

% From Math 55 and Math 145 at Harvard
\newenvironment{subproof}[1][Proof]{%
\begin{proof}[#1] \renewcommand{\qedsymbol}{$\blacksquare$}}%
{\end{proof}}

\newcommand{\liff}{\leftrightarrow}
\newcommand{\lthen}{\rightarrow}
\newcommand{\opname}{\operatorname}
\newcommand{\surjto}{\twoheadrightarrow}
\newcommand{\injto}{\hookrightarrow}
\newcommand{\On}{\mathrm{On}} % ordinals
\DeclareMathOperator{\img}{im} % Image
\DeclareMathOperator{\Img}{Im} % Image
\DeclareMathOperator{\coker}{coker} % Cokernel
\DeclareMathOperator{\Coker}{Coker} % Cokernel
\DeclareMathOperator{\Ker}{Ker} % Kernel
\DeclareMathOperator{\rank}{rank}
\DeclareMathOperator{\Spec}{Spec} % spectrum
\DeclareMathOperator{\Tr}{Tr} % trace
\DeclareMathOperator{\pr}{pr} % projection
\DeclareMathOperator{\ext}{ext} % extension
\DeclareMathOperator{\pred}{pred} % predecessor
\DeclareMathOperator{\dom}{dom} % domain
\DeclareMathOperator{\ran}{ran} % range
\DeclareMathOperator{\Hom}{Hom} % homomorphism
\DeclareMathOperator{\End}{End} % endomorphism

% Things Lie
\newcommand{\kb}{\mathfrak b}
\newcommand{\kg}{\mathfrak g}
\newcommand{\kh}{\mathfrak h}
\newcommand{\kn}{\mathfrak n}
\newcommand{\ku}{\mathfrak u}
\newcommand{\kz}{\mathfrak z}
\DeclareMathOperator{\Ext}{Ext} % Ext functor
\DeclareMathOperator{\Tor}{Tor} % Tor functor
\newcommand{\gl}{\opname{\mathfrak{gl}}} % frak gl group
\renewcommand{\sl}{\opname{\mathfrak{sl}}} % frak gl group

% More script letters etc.
\newcommand{\SA}{\mathscr A}
\newcommand{\SB}{\mathscr B}
\newcommand{\SC}{\mathscr C}
\newcommand{\SD}{\mathscr D}
\newcommand{\SF}{\mathscr F}
\newcommand{\SG}{\mathscr G}
\newcommand{\SH}{\mathscr H}
\newcommand{\OO}{\mathcal O}

%% Napkin commands
\newcommand{\prototype}[1]{
	\emph{{\color{red} Prototypical example for this section:} #1} \par\medskip
}
\newenvironment{moral}{%
	\begin{mdframed}[linecolor=green!70!black]%
	\bfseries\color{green!70!black}}%
	{\end{mdframed}}

%%% END EXCERPT OF EVAN.STY
%%fakesection Commutative diagrams
\usepackage{diagrams}
\newarrow{Inj}C--->
\newarrow{Surj}----{>>}
\newarrow{Proj}----{>>}
\newarrow{Embed}>--->
\newarrow{Incl}C--->
\newarrow{Mapsto}|--->
\newarrow{Isom}<--->
\newarrow{Dotted}....>
\newarrow{Dashed}{}{dash}{}{dash}>
\newarrow{Line}-----
\usepackage{tikz-cd}
\usetikzlibrary{decorations.pathmorphing}
\tikzcdset{row sep/normal=3.14em,column sep/normal=3.14em}

%%fakesection Page layout
\usepackage[headsepline]{scrlayer-scrpage}
\renewcommand{\headfont}{}
\addtolength{\textheight}{3.14cm}
\setlength{\footskip}{0.5in}
\setlength{\headsep}{10pt}
\ihead{Evan Chen}
\automark{section}
\chead{}
\ohead{\headmark}
\cfoot{\pagemark}

%%fakesection Fancy section and chapter heads
\addtokomafont{chapterprefix}{\raggedleft}
\renewcommand*{\sectionformat}{\color{purple}\S\thesection\autodot\enskip}

\RedeclareSectionCommand[beforeskip=0.5em]{chapter}
\renewcommand*{\chapterformat}{%
\mbox{\scalebox{1.5}{\chapappifchapterprefix{\nobreakspace}}%
\scalebox{2.718}{\color{purple}\thechapter\autodot}\enskip}}

\addtokomafont{partprefix}{\rmfamily}
\renewcommand*{\partformat}{\color{purple}\scalebox{2.5}{\thepart}%
	\enlargethispage{2em}}
\newcommand{\problemhead}{Problems to think about}

%%fakesection Mini ToC
\usepackage[tight]{minitoc}
\noptcrule
\mtcsetfont{parttoc}{section}{\small\rmfamily\upshape\mdseries}
\renewcommand*{\partheadstartvskip}{}
\renewcommand*{\partheadendvskip}{}
\renewcommand\beforeparttoc{\hspace{\fill}\rule{0.95\linewidth}{2pt}\hspace{\fill}}
\let\oldpart\part
\renewcommand{\part}[1]{\oldpart{#1}\parttoc\cleardoublepage}

%%fakesection SNSD
\newcommand{\snsd}[2][]%
{\includegraphics[#1]{/home/evan/Pictures/TopologicalGG/#2}}

%%fakesection Theorems
\usepackage{thmtools}
\usepackage[framemethod=TikZ]{mdframed}

\theoremstyle{definition}
\mdfdefinestyle{mdbluebox}{%
	roundcorner = 10pt,
	linewidth=1pt,
	skipabove=12pt,
	innerbottommargin=9pt,
	skipbelow=2pt,
	nobreak=true,
	linecolor=blue,
	backgroundcolor=TealBlue!5,
}
\declaretheoremstyle[
	headfont=\sffamily\bfseries\color{MidnightBlue},
	mdframed={style=mdbluebox},
	headpunct={\\[3pt]},
	postheadspace={0pt}
]{thmbluebox}

\mdfdefinestyle{mdrecbox}{%
	linewidth=0.5pt,
	skipabove=12pt,
	frametitleaboveskip=5pt,
	frametitlebelowskip=0pt,
	skipbelow=2pt,
	frametitlefont=\bfseries,
	innertopmargin=4pt,
	innerbottommargin=8pt,
	nobreak=true,
	linecolor=RawSienna,
	backgroundcolor=Salmon!5,
}
\declaretheoremstyle[
	headfont=\bfseries\color{RawSienna},
	mdframed={style=mdrecbox},
	headpunct={\\[3pt]},
	postheadspace={0pt},
]{thmrecbox}

\declaretheorem[%
style=thmbluebox,name=Theorem,numberwithin=section]{theorem}
\declaretheorem[style=thmbluebox,name=Lemma,sibling=theorem]{lemma}
\declaretheorem[style=thmbluebox,name=Proposition,sibling=theorem]{proposition}
\declaretheorem[style=thmbluebox,name=Corollary,sibling=theorem]{corollary}
\declaretheorem[style=thmrecbox,name=Example,sibling=theorem]{example}

\mdfdefinestyle{mdinlinebox}{%
	skipabove=8pt,
	linewidth=3pt,
	rightline=false,
	leftline=true,
	topline=false,
	bottomline=false,
	linecolor=gray,
	backgroundcolor=black!2
}
\declaretheoremstyle[
	headfont=\bfseries,
	bodyfont=\normalfont\small,
	spaceabove=0pt,
	spacebelow=0pt,
	mdframed={style=mdinlinebox}
]{thminlinebox}
\theoremstyle{theorem}
\declaretheorem[name=Question,sibling=theorem,style=thminlinebox]{ques}
\declaretheorem[name=Exercise,sibling=theorem,style=thminlinebox]{exercise}

\theoremstyle{definition}
\newtheorem{claim}[theorem]{Claim}
\newtheorem{definition}[theorem]{Definition}
\newtheorem{fact}[theorem]{Fact}
\newtheorem{remark}[theorem]{Remark}
\newtheorem{abuse}[theorem]{Abuse of Notation}

\newtheorem{problem}{Problem}[chapter]
\renewcommand{\theproblem}{\thechapter\Alph{problem}}
\newtheorem{sproblem}[problem]{Problem}
\newtheorem{dproblem}[problem]{Problem}
\renewcommand{\thesproblem}{\theproblem$^{\star}$}
\renewcommand{\thedproblem}{\theproblem$^{\dagger}$}
\newcommand{\listhack}{$\empty$\vspace{-2em}}

%%fakesection Answers
\usepackage{answers}
\Newassociation{hint}{answeritem}{tex/backmatter/all-hints}
\Newassociation{sol}{answeritem}{tex/backmatter/all-solns}
\renewcommand{\solutionextension}{out}
\renewenvironment{answeritem}[1]{\item[\bf #1.]}{}

%%fakesection Hot chili peppers
\reversemarginpar
\newcommand{\prechili}{\vspace*{0.3em}\hspace*{1.5em}}
\newcommand{\nochili}{\hspace*{1.5em}}
\newcommand{\chili}{\includegraphics[width=1.5em]{media/chili.png}}
\newcommand{\gim}{\marginpar{\prechili\nochili\nochili\chili}}
\newcommand{\yod}{\marginpar{\prechili\nochili\chili\chili}}
\newcommand{\kurumi}{\marginpar{\prechili\chili\chili\chili}}

%%fakesection Table of contents
% First add ToC to ToC
\makeatletter
\usepackage{etoolbox}
\pretocmd{\tableofcontents}{%
  \if@openright\cleardoublepage\else\clearpage\fi
  \pdfbookmark[0]{\contentsname}{toc}%
}{}{}%
\makeatother
\setcounter{tocdepth}{1}
% Fix spacing issues in To
\usepackage[tocindentauto]{tocstyle}
\usetocstyle{KOMAlike}

%%fakesection Asymptote definitions
\usepackage{patch-asy}
\numberwithin{asy}{chapter}
\renewcommand{\theasy}{\thechapter\Alph{asy}}
\begin{asydef}
	import olympiad;
	import cse5;
	size(6cm);
	usepackage("amsmath");
	usepackage("amssymb");
	defaultpen(fontsize(11pt));
	pointfontsize=11;
	void bigbox(string L) {
		filldraw((-4,3)--(4,3)--(4,-3)--(-4,-3)--cycle,
			lightblue+opacity(0.2), black);
		label(L, (4,3), dir(-45));
	}
	void bigblob(string L) {
		filldraw((-4,3)..(-5,1)..(-4.7,0)..(-4.1,-2)..(-4,-3)
			..(-1,-3.3)..(1,-3.2)..(4,-3)
			..(4.3,-1)..(4.1,0.5)..(4,3)
			..(2,3.3)..(0,3.1)..(-3,2.9)..cycle,
			lightblue+opacity(0.2), black);
		label(L, (4,3), dir(5));
	}

	real[] default_ticks = {-4,-3,-2,-1,0,1,2.3,4};
	void cplane(real theta=-45, real[] T = default_ticks, real xmin = -4.5, real xmax = 4.5, real S = 3) {
		graph.xaxis("Re", xmin, xmax, graph.Ticks(format="%", Ticks=T, Size=S), Arrows);
		graph.yaxis("Im", xmin, xmax, graph.Ticks(format="%", Ticks=T, Size=S), Arrows);
		Drawing("0", origin, dir(theta));
		markscalefactor *= xmax;
	}

	pair opendot(pair P, pen p = defaultpen) {
		dot(P, p, UnFill);
		dot(P, p, Draw);
		return P;
	}
\end{asydef}
\def\asydir{asy}

%%fakesection Bibliography
\usepackage[backend=biber,style=alphabetic]{biblatex}
\DeclareLabelalphaTemplate{
	\labelelement{
		\field[final]{shorthand}
		\field{label}
		\field[strwidth=2,strside=left]{labelname}
	}
	\labelelement{
		\field[strwidth=2,strside=right]{year}    
	}
}
\DeclareFieldFormat{labelalpha}{\textbf{\scriptsize #1}}
\addbibresource{references.bib}
\addbibresource{images.bib}

%%fakesection Napkin Macros
\newcommand{\Zc}[1]{\ZZ_{#1}}
\newcommand{\Zm}[1]{\ZZ_{#1}^\times}
\newcommand{\pre}{^{\text{pre}}}
\newcommand{\normalin}{\trianglelefteq}
\newcommand{\triv}{\mathrm{triv}}
\newcommand{\largeotimes}[1]{\mathop{\otimes}\limits_{#1}}
\newcommand{\Mat}{\mathrm{Mat}}
\newcommand{\PGL}[1]{\mathbf{PGL}_{#1}(\CC)}

\DeclareMathOperator{\Stab}{Stab}
\DeclareMathOperator{\FixPt}{FixPt}
\DeclareMathOperator{\refl}{refl}
\DeclareMathOperator{\Fun}{Fun}
\DeclareMathOperator{\Irrep}{Irrep}
\DeclareMathOperator{\Res}{Res}
\DeclareMathOperator{\Reg}{Reg}
\DeclareMathOperator{\Classes}{Classes}
\newcommand{\FunCl}{\Fun_{\mathrm{class}}}
\newcommand{\Homrep}{\Hom_{\mathrm{rep}}}
\newcommand{\ab}{^\text{ab}} % abelianization
\DeclareMathOperator{\ev}{ev}
\DeclareMathOperator{\Wind}{\mathbf I}
\newcommand{\Ctriv}{\CC_{\mathrm{triv}}}
\newcommand{\Csign}{\CC_{\mathrm{sign}}}

% Alg geom macros
\newcommand{\VV}{\mathcal V}
\newcommand{\Vp}{\mathcal V_{\text{pr}}}
\newcommand{\Aff}{\mathbb A}
\newcommand{\II}{\mathcal I_{\text{rad}}}
\newcommand{\RP}{\mathbb{RP}}
\newcommand{\CP}{\mathbb{CP}}
\newcommand{\Cl}{\mathrm{Cl}}
\newcommand{\restrict}[1]{\bgroup\restriction_{#1}\egroup}
\DeclareMathOperator{\Opens}{OpenSets}
\DeclareMathOperator{\Proj}{Proj}
\DeclareMathOperator{\res}{res}
\newcommand{\mm}{\mathfrak m}
\newcommand{\sh}{^\mathrm{sh}}

\renewcommand{\Re}{\opname{Re}}
\renewcommand{\Im}{\opname{Im}}

%% Category Theory macros
\renewcommand{\AA}{\mathcal A}
\newcommand{\BB}{\mathcal B}
\newcommand{\obj}{\operatorname{obj}}
\newcommand{\op}{^{\mathrm{op}}}


\usepackage{tkz-graph}
\pgfarrowsdeclare{biggertip}{biggertip}{%
  \setlength{\arrowsize}{1pt}  
  \addtolength{\arrowsize}{.1\pgflinewidth}  
  \pgfarrowsrightextend{0}  
  \pgfarrowsleftextend{-5\arrowsize}  
}{%
  \setlength{\arrowsize}{1pt}  
  \addtolength{\arrowsize}{.1\pgflinewidth}  
  \pgfpathmoveto{\pgfpoint{-5\arrowsize}{4\arrowsize}}  
  \pgfpathlineto{\pgfpointorigin}  
  \pgfpathlineto{\pgfpoint{-5\arrowsize}{-4\arrowsize}}  
  \pgfusepathqstroke  
}  
\tikzset{
	EdgeStyle/.style = {>=biggertip}
}

\usepackage[all,cmtip,2cell]{xy}
\UseTwocells
\newcommand\nattfm[5]{\xymatrix@C+2pc{#1 \rtwocell<4>^{#2}_{#4}{\; #3} & #5}}

%% Alg Top macros
\DeclareMathOperator{\Cells}{Cells}
\newcommand{\HdR}{H_{\mathrm{dR}}}

%% Alg NT macros
\renewcommand{\aa}{\mathfrak a}
\newcommand{\bb}{\mathfrak b}
\newcommand{\pp}{\mathfrak p}
\newcommand{\qq}{\mathfrak q}

\newcommand{\Frob}{\mathrm{Frob}}
\DeclareMathOperator{\Norm}{N}
\DeclareMathOperator{\Ram}{Ram}
\newcommand{\TrK}{\Tr_{K/\QQ}}
\newcommand{\NK}{\Norm_{K/\QQ}}
\newcommand{\kf}{\mathfrak f}
\newcommand{\kP}{\mathfrak P}
\newcommand{\kQ}{\mathfrak Q}

%% Diff geo macros
\newcommand{\ee}{\mathbf e}
\newcommand{\fpartial}[2]{\frac{\partial #1}{\partial #2}}

%% ST macros
\newcommand{\CH}{\mathsf{CH}}
\newcommand{\ZFC}{\mathsf{ZFC}}

\newcommand{\Name}{\text{Name}}
\newcommand{\Po}{\mathbb P}
\newcommand{\nrank}{\opname{n-rank}} % ranks
\newcommand{\PP}{\mathcal P}

\newcommand{\EmptySet}{\mathrm{EmptySet}}
\newcommand{\PowerSet}{\mathrm{PowerSet}}
\newcommand{\Pairing}{\mathrm{Pairing}}
\newcommand{\Infinity}{\mathrm{Infinity}}
\newcommand{\Extensionality}{\mathrm{Extensionality}}
\newcommand{\Foundation}{\mathrm{Foundation}}
\newcommand{\Union}{\mathrm{Union}}
\newcommand{\Comprehension}{\mathrm{Comprehension}}
\newcommand{\Replacement}{\mathrm{Replacement}}

\usepackage{mathrsfs}

\newcommand\MM{\mathscr M}
\newcommand\llex{<_{\text{lex}}}

\DeclareMathOperator{\cof}{cof}

%% Quantum macros
\usepackage{braket}
\newcommand{\cvec}[1]{\begin{bmatrix} #1 \end{bmatrix}}
\newcommand{\pair}[2]{\begin{bmatrix} #1 \\ #2 \end{bmatrix}}
\newcommand{\zup}{\ket\uparrow}
\newcommand{\zdown}{\ket\downarrow}
\newcommand{\xup}{\ket\rightarrow}
\newcommand{\xdown}{\ket\leftarrow}
\newcommand{\yup}{\ket\otimes}
\newcommand{\ydown}{\ket\odot}
\newcommand{\UCNOT}{U_{\mathrm{CNOT}}}
\newcommand{\UQFT}{U_{\mathrm{QFT}}}

\def\asydir{}
\addbibresource{../../references.bib}
\renewcommand{\gim}{($\ast$)}


\begin{document}
\title{Quantum Algorithms}
\maketitle

\chapter{Quantum States and Measurements}
In this chapter we'll explain how to set up quantum states using
linear algebra. This will allow me to talk about quantum \emph{circuits}
in the next chapter, which will set the stage for Shor's algorithm.

I won't do very much physics (read: none at all).
That is, I'll only state what the physical reality is in terms
of linear algebras, and defer the philosophy of why this is true
to your neighborhood ``Philosophy of Quantum Mechanics'' class
(which is a ``social science'' class at MIT!).

\section{Bra-Ket Notation}
Physicists have their own notation for vectors:
whereas I previously used something like $v$, $e_1$, and so on,
in this chapter you'll see the infamous \vocab{bra-ket} notation:
a vector will be denoted by $\ket\bullet$, where $\bullet$ is some
variable name: unlike in math or Python, this can include
numbers, symbols, Unicode characters, whatever you like.
This is called a ``ket''.
To pay a homage to physicists everywhere,
we'll use this notation for this chapter too.

For us:
\begin{definition}
	A \vocab{Hilbert space} $H$ is a finite-dimensional
	complex inner product space.
\end{definition}
\begin{abuse}
	In actual physics, almost all Hilbert spaces used
	are infinite-dimensional, in which case there is 
	the additional requirement that $H$ be a complete metric
	space with respect to its norm. But finite-dimensional spaces
	will suffice for the purposes here
	(and automatically satisfy the completeness property).
\end{abuse}

If $\dim H = n$, then its orthonormal basis elements are often denoted
\[ \ket0, \ket1, \dots, \ket{n-1} \]
(instead of $e_i$)
and a generic element of $H$ denoted by
\[ \ket\psi, \ket\phi, \dots \]
and various other Greek letters.

Now for any $\ket\psi \in H$,
we can consider the canonical dual element in $H^\vee$
(since $H$ has an inner form), which we denote by $\bra\psi$ (a ``bra'').
For example, if $\dim H = 2$ then we can write
\[ \ket\psi = \begin{pmatrix} \alpha \\ \beta \end{pmatrix} \]
in an orthonormal basis, in which case
\[ \bra\psi = \begin{pmatrix} \ol\alpha & \ol\beta \end{pmatrix}. \]
We even can write dot products succinctly in this notation:
if $\ket\phi = \begin{pmatrix} \gamma \\ \delta \end{pmatrix}$,
then the dot product of $\ket\phi$ and $\ket\psi$ is given by
\[
	\braket{\psi|\phi}
	= \cvec{\ol\alpha & \ol\beta} \cvec{\gamma \\ \delta}
	= \ol\alpha\gamma + \ol\beta \delta.
\]
So we will use the notation $\braket{\psi|\phi}$
instead of the more mathematical $\left< \ket\psi, \ket\phi \right>$.
In particular, the squared norm of $\ket\psi$ is just $\braket{\psi|\psi}$.
Concretely, for $\dim H = 2$ we have
$\braket{\psi|\psi} = |\alpha^2| + |\beta|^2$.

\section{The State Space}
If you think that's weird, well, it gets worse.

In classical computation, a bit is either $0$ or $1$.
More generally, we can think of a classical space of $n$
possible states $0$, \dots, $n-1$.
Thus in the classical situation, the space of possible states
is just a discrete set with $n$ elements.

In quantum computation, a \vocab{qubit} is instead
any \emph{complex linear combination} of $0$ and $1$.
To be precise, consider the normed complex vector space
\[ H = \CC^{\oplus 2} \]
and denote the orthonormal basis elements by $\ket0$ and $\ket1$.
Then a \emph{qubit} is a nonzero element $\ket\psi \in H$,
so that it can be written in the form
\[ \ket\psi = \alpha \ket 0 + \beta \ket 1 \]
where $\alpha$ and $\beta$ are not both zero.
Typically, we normalize so that $\ket\psi$ has norm $1$:
\[ \braket{\psi|\psi} = 1 \iff |\alpha|^2 + |\beta|^2 = 1. \]
In particular, we can recover the ``classical'' situation
with $\ket 0 \in H$ and $\ket 1 \in H$,
but now we have some ``intermediate'' states,
such as \[ \frac{1}{\sqrt2} \left(\ket 0 + \ket 1 \right). \]
Philosophically, what has happened is that:
\begin{moral}
	Instead of allowing just the states $\ket 0$ and $\ket 1$,
	we allow any complex linear combination of them.
\end{moral}
More generally, if $\dim H = n$,
then the possible states are nonzero elements
\[ c_0\ket0 + c_1\ket1 + \dots + c_{n-1}\ket{n-1} \]
which we usually normalize so that
$|c_0|^2 + |c_1|^2 + \dots + |c_{n-1}|^2 = 1$.

\section{Observations}
\prototype{$\id$ corresponds to not making a measurement
	since all its eigenvalues are equal,
	but any operator with distinct eigenvalues will cause collapse.}
If you think that's weird, well, it gets worse.
First, some linear algebra:
\begin{definition}
	Let $V$ be a finite-dimensional normed vector space.
	For a map $T: V \to V$, the following conditions are equivalent:
	\begin{itemize}
		\ii $\left< Tx, y\right> = \left< x, Ty \right>$
		for any $x,y \in V$.
		\ii $T = T^\dagger$.
	\end{itemize}
	A map $T$ satisfying these conditions is called \vocab{Hermitian}.
\end{definition}
\begin{ques}
	Show that $T$ is normal.
\end{ques}
Thus, we know that $T$ is diagonalizable
with respect to the inner form, so for a suitable basis we
can write it in an orthonormal basis as
\[
	T = \begin{pmatrix}
		\lambda_0 & 0 & \dots & 0 \\
		0 & \lambda_1 & \dots & 0 \\
		\vdots & \vdots & \ddots & \vdots \\
		0 & 0 & \dots & \lambda_{n-1}
	\end{pmatrix}.
\]
As we've said, this is fantastic:
not only do we have a basis of eigenvectors,
but the eignvectors are pairwise orthogonal,
and so they form an orthonormal basis of $V$.
\begin{ques}
	Show that all eigenvalues of $T$ are real.
	($T = T^\dagger$.)
\end{ques}

Back to quantum computation.
Suppose we have a state $\ket\psi \in H$, where $\dim H = 2$;
we haven't distinguished a particular basis yet,
so we just have a nonzero vector.
Then the way observations work (and this is physics, so you'll have to
take my word for it) is as follows:
\begin{moral}
	Pick a Hermitian operator $T : H \to H$;
	then observations of $T$ return eigenvalues of $T$.
\end{moral}
To be precise:
\begin{itemize}
	\ii Pick a Hermitian operator $T : H \to H$,
	which is called the \vocab{observable}.
	\ii Consider its eigenvalues $\lambda_0$, \dots, $\lambda_{n-1}$
	and corresponding eigenvectors $\ket{0}_T$, \dots, $\ket{n-1}_T$.
	Tacitly we may assume that $\ket{0}_T$, \dots, $\ket{n-1}_T$ form
	an orthonormal basis of $H$.
	(The subscript $T$ is here to distinguish the eigenvectors of $T$
	from the basis elements of $H$.)
	\ii Write $\ket\psi$ in the orthonormal basis as
	\[ c_0\ket0_T + c_1\ket1_T + \dots + c_{n-1}\ket{n-1}_T. \]
	\ii Then the probability of observing $\lambda_i$ is
	\[ \frac{|c_i|^2}{|c_0|^2 + \dots + |c_{n-1}|^2}. \]
	This is called making an \vocab{observation along $T$}.
\end{itemize}
Note that in particular, for any nonzero constant $c$,
$\ket\psi$ and $c\ket\psi$ are indistinguishable,
which is why we like to normalize $\ket\psi$.
But the queerest thing of all is what happens to $\ket\psi$:
by measuring it, we actually destroy information.
This behavior is called \vocab{quantum collapse}.
\begin{itemize}
	\ii Suppose for simplicity that we observe $\ket\psi$
	with $T$ and obtain an eigenvalue $\lambda$,
	and that $\ket{i}_T$ is the only eigenvector with this eigenvalue.
	Then, the state $\ket\psi$ \emph{collapses} to just the state
	$c_i \ket{i}_T$: all the other information is destroyed.
	(In fact, we may as well say it collapses to $\ket{i}_T$,
	since again constant factors are not relevant.)

	\ii	More generally, if we observe $\lambda$,
	consider the generalized eigenspace $H_\lambda$, i.e.\
	(i.e.\ the span of eigenvectors with the same eigenvalue).
	Then the physical state $\ket\psi$ has been changed as well:
	it has now been projected onto the eigenspace $H_\lambda$.
	In still other words, after observation, the state collapses to
	\[
		\sum_{\substack{0 \le i \le n \\ \substack{\lambda_i = \lambda}}}
		c_i \ket{i}_T.
	\]
	All the coefficients from different eigenspaces have been destroyed.
\end{itemize}
\begin{ques}
	Suppose we measure $\ket\psi$ with $T$ and get $\lambda$.
	What happens if we measure with $T$ again?
\end{ques}
Why does this happen? Beats me\dots physics (and hence real life) is weird.
But anyways, an example.
\begin{example}
	[Quantum Measurement of a state $\ket\psi$]
	Let $H = \CC^{\oplus 2}$ with orthonormal basis $\ket0$ and $\ket1$
	and consider the state
	\[
		\ket\psi
		= \frac{i}{\sqrt5} \ket 0
		+ \frac{2}{\sqrt5} \ket 1
		= \pair{i/\sqrt5}{2/\sqrt5} \in H.
	\]
	\begin{enumerate}[(a)]
		\ii Let \[ T = \begin{pmatrix} 1 & 0 \\ 0 & -1 \end{pmatrix}. \]
		This has eigenvectors $\ket0 = \ket0_T$ and $\ket1 = \ket1_T$,
		with eigenvalues $+1$ and $-1$.  So if we measure $\ket\psi$ to $T$,
		we get $+1$ with probability $1/5$ and $-1$ with probability $4/5$.
		After this measurement, the original state collapses to
		$\ket0$ if we measured $0$, and $\ket1$ if we measured $1$.
		So we never learn the original probabilities.

		\ii Now consider $T = \id$, and arbitrarily
		pick two orthonormal eigenvectors $\ket0_T$, $\ket1_T$;
		thus $\psi = c_0\ket0_T + c_1\ket1_T$.
		Since all eigenvalues of $T$ are $+1$,
		our measurement will always be $+1$ no matter what we do.
		But there is also no collapsing,
		because none of the coefficients get destroyed.

		\ii Now consider
		\[ T = \begin{pmatrix} 0 & 7 \\ 7 & 0 \end{pmatrix}. \]
		The two normalized eigenvectors are
		\[ \ket0_T = \frac{1}{\sqrt2}\pair11
		\qquad \ket1_T = \frac{1}{\sqrt2}\pair1{-1} \]
		with eigenvalues $+7$ and $-7$ respectively. In this basis, we have
		\[
			\ket\psi = \frac{2+i}{\sqrt{10}}\ket0_T
			+ \frac{-2+i}{\sqrt{10}}\ket1_T. \]
		So we get $+7$ with probability $\half$ and $-7$
		with probability $\half$, and after the measurement,
		$\ket\psi$ collapses to one of $\ket0_T$ and $\ket1_T$.
	\end{enumerate}
\end{example}

For $H = \CC^{\oplus 2}$ we can come up with more classes of examples using
the so-called \vocab{Pauli matrices}.
These are the three Hermitian matrices
\[
	\sigma_z = \begin{pmatrix} 1 & 0 \\ 0 & -1 \end{pmatrix}
	\qquad
	\sigma_x = \begin{pmatrix} 0 & 1 \\ 1 & 0 \end{pmatrix}
	\qquad
	\sigma_y = \begin{pmatrix} 0 & i \\ -i & 0 \end{pmatrix}.
\]
\begin{ques}
	Show that these three matrices, plus the identity matrix,
	form a basis for the set of Hermitian $2 \times 2$ matrices.
\end{ques}
Their normalized eigenvectors are
\[ \zup = \ket0 = \pair10 \qquad \zdown = \ket1 = \pair01 \]
\[ \xup = \frac{1}{\sqrt2}\pair11
	\qquad \xdown = \frac{1}{\sqrt2}\pair1{-1} \]
\[ \yup = \frac{1}{\sqrt2}\pair1i
	\qquad \ydown = \frac{1}{\sqrt2}\pair1{-i} \]
which we call ``$z$-up'', ``$z$-down'',
``$x$-up'', ``$x$-down'', ``$y$-up'', ``$y$-down''.
So, given a state $\ket\psi \in \CC^{\oplus 2}$
we can make a measurement with respect to any of these three bases
by using the corresponding Pauli matrix.

In light of this, the previous examples were (a) measuring along $\sigma_z$,
(b) measuring along $\id$, and (c) measuring along $7\sigma_x$.

Notice that if we are given a state $\ket\psi$,
and are told in advance that it is either $\xup$ or $\xdown$,
for example, then we are in what is more or less a classical situation.
Specifically, if we make a measurement along $\sigma_x$,
then we find out which state that $\ket\psi$ was in (with 100\% certainty),
and the state does not undergo any collapse.

\section{Entanglement}
\prototype{Singlet state: spooky action at a distance.}
If you think that's weird, well, it gets worse.

Qubits don't just act independently:
they can talk to each other by means of a \emph{tensor product}.
Explicitly, consider \[ H = \CC^{\oplus 2} \otimes \CC^{\oplus 2} \]
endowed with the norm described in \Cref{prob:inner_prod_tensor}.
One should think of this as a qubit $A$ in a space $H_A$
along with a second qubit $B$ in a different space $H_B$,
which have been allowed to interact in some way,
and $H = H_A \otimes H_B$ is the set of possible states of \emph{both} qubits.
Thus
\[ 
	\ket{0}_A \otimes \ket{0}_B, \quad
	\ket{0}_A \otimes \ket{1}_B, \quad
	\ket{1}_A \otimes \ket{0}_B, \quad
	\ket{1}_A \otimes \ket{1}_B
\]
is an orthonormal basis of $H$;
here $\ket{i}_A$ is the basis of the first $\CC^{\oplus 2}$
while $\ket{i}_B$ is the basis of the second $\CC^{\oplus 2}$,
so these vectors should be thought of as ``unrelated''
just as with any tensor product.

As before, a measurement of a state in $H$ requires
a Hermitian map $H \to H$.
In particular, if we only want to measure the qubit $B$ along $M_B$,
we can use the operator \[ \id_A \otimes M_B. \]
The eigenvalues of this operator coincide with the ones for $M_B$,
and the eigenspace for $\lambda$ will be the $H_A \otimes (H_B)_\lambda$,
so when we take the projection the $A$ qubit will be unaffected.

This does what you would hope for pure tensors in $H$:
\begin{example}[Two Non-Entangled Qubits]
	Suppose we have qubit $A$ in the state
	$\frac{i}{\sqrt5}\ket0_A + \frac{2}{\sqrt5}\ket1_A$
	and qubit $B$ in the state
	$\frac{1}{\sqrt2} \ket0_B + \frac{1}{\sqrt2}\ket1_B$.
	So, the two qubits in tandem are represented by the tensor product
	\[
		\ket\psi
		=
		\left( \frac{i}{\sqrt5}\ket0_A + \frac{2}{\sqrt5}\ket1_A \right)
		\otimes
		\left( \frac{1}{\sqrt2} \ket0_B + \frac{1}{\sqrt2}\ket1_B \right).
	\]
	Suppose we measure $\ket\psi$ along 
	\[ M = \id_A \otimes \sigma_z^B. \]
	The eigenspace decomposition is
	\begin{itemize}
		\ii $+1$ for the span of $\ket0_A \otimes \ket0_B$ and
		$\ket1_A \otimes \ket0_B$, and
		\ii $-1$ for the span of $\ket0_A \otimes \ket1_B$ and
		$\ket1_A \otimes \ket1_B$.
	\end{itemize}
	(We could have used other bases, like $\xup_A \otimes \ket0_B$ and
	$\xdown_A \otimes \ket0_B$ for the first eigenpsace, but no worries.)
	Expanding $\ket\psi$ in the four-element basis, we find that
	we'll get the first eigenspace with probability
	\[ \left|\frac{i}{\sqrt{10}}\right|^2
	+ \left|\frac{2}{\sqrt{10}}\right|^2 = \half. \]
	and the second eigenspace with probability $\half$ as well.
	(Note how the coefficients for $A$ don't do anything!)
	After the measurement, we destroy the coefficients of the other eigenspace;
	thus (after re-normalization) we obtain the collapsed state
	\[ \left( \frac{i}{\sqrt5}\ket0_A + \frac{2}{\sqrt5}\ket1_A \right)
		\otimes \ket0_B
		\qquad\text{or}\qquad
		\left( \frac{i}{\sqrt5}\ket0_A + \frac{2}{\sqrt5}\ket1_A \right)
		\otimes \ket1_B
	\]
	again with 50\% probability each.
\end{example}
So this model lets us more or less work with the two qubits independently:
when we make the measurement, we just make sure to not touch the other qubit
(which corresponds to the identity operator).

\begin{exercise}
	Show that if $\id_A \otimes \sigma_x^B$ is applied to the $\ket\psi$
	in this example, there is no collapse at all.
	What's the result of this measurement?
\end{exercise}

However, some states behave very weirdly.
Consider the state
\[
	\ket{\Psi_-}
	=
	\frac{1}{\sqrt2} \ket0_A \otimes \ket1_B
	- \frac{1}{\sqrt2} \ket1_A \otimes \ket0_B
\]
which is called the \vocab{singlet state}.
One can see that $\ket{\Psi_-}$ is not a simple tensor,
which means that it doesn't just consists of two qubits side by side:
the qubits in $H_A$ and $H_B$ have become \emph{entangled}.

Now, what happens if we measure just the qubit $A$?
This corresponds to making the measurement
\[ T = \sigma_z^A \otimes \id_B. \]
The eigenspace decomposition of $T$ can be described as:
\begin{itemize}
	\ii The space spanned by $\ket{0}_A \otimes \ket{0}_B$ and
	$\ket{0}_A \otimes \ket{1}_B$, with eigenvalue $+1$.
	\ii The space spanned by $\ket{1}_A \otimes \ket{0}_B$ and
	$\ket{1}_A \otimes \ket{1}_B$, with eigenvalue $-1$.
\end{itemize}
So one of two things will happen:
\begin{itemize}
	\ii With probability $\half$, we measure $+1$
	and the collapsed state is $\ket0_A \otimes \ket1_B$.
	\ii With probability $\half$, we measure $-1$
	and the collapsed state is $\ket1_A \otimes \ket0_B$.
\end{itemize}
But now we see that measurement along $A$ has told us what the
state of the bit $B$ is completely!
By solely looking at measurements on $A$, we learn $B$;
this paradox is called \emph{spooky action at a distance},
or in Einstein's tongue, \vocab{spukhafte Fernwirkung}.  

Thus,
\begin{moral}
	In tensor products of Hilbert spaces,
	states which are not pure tensors
	correspond to ``entangled'' states.
\end{moral}

What this really means is that the qubits cannot be described independently;
the state of the system must be given as a whole.

\section\problemhead

\begin{problem}
	We measure $\ket{\Psi_-}$ by $\sigma_x^A \otimes \id_B$,
	and hence obtain either $+1$ or $-1$.
	Show that the state of qubit $B$ is determined
	by this measurement completely, and find the corresponding values.
	\begin{hint}
		Rewrite $\ket{\Psi_-} = -\frac{1}{\sqrt2}
		\left( \xup_A\otimes\xdown_B - \xdown_A\xup_B \right)$.
	\end{hint}
	\begin{sol}
		By a straightforward computation, 
		we have $\ket{\Psi_-} = -\frac{1}{\sqrt2}
		\left( \xup_A\otimes\xdown_B - \xdown_A\xup_B \right)$.
		Now, $\xup_A \otimes \xup_B$, $\xup_A \otimes \xdown_B$
		span one eigenspace of $\sigma_x^A \otimes \id_B$,
		and $\xdown_A \otimes \xup_B$, $\xdown_A \otimes \xdown_B$
		span the other. So this is the same as before.
	\end{sol}
\end{problem}

\begin{problem}
	[Greenberger-Horne-Zeilinger Paradox]
	Consider the state in $(\CC^{\oplus2})^{\otimes 3}$
	\[
		\ket{\Psi}_{\text{GHZ}}
		=
		\frac{1}{\sqrt2}
		\left( 
		\ket0_A \ket0_B \ket0_C
		- \ket1_A \ket1_B \ket1_C \right).
	\]
	Find the value of the measurements along each of
	\[ \sigma_y^A \otimes \sigma_y^B \otimes \sigma_x^C , \quad
		\sigma_y^A \otimes \sigma_x^B \otimes \sigma_y^C, \quad
		\sigma_x^A \otimes \sigma_y^B \otimes \sigma_y^C, \quad
		\sigma_x^A \otimes \sigma_x^B \otimes \sigma_x^C.
	\]
\end{problem}

\chapter{Quantum Circuits}
\chapter{Shor's Algorithm}

\end{document}
